\subsection*{Serverless Computing e FaaS nel continuum Edge--Cloud}

Parallelamente all'evoluzione verso l'edge, il paradigma
\emph{Serverless Computing} ha acquisito crescente rilevanza
come modello di programmazione cloud-native.
Nel modello serverless, lo sviluppatore distribuisce codice
senza occuparsi della gestione dell'infrastruttura sottostante,
che viene completamente delegata al provider
\cite{VanEyk2018_ServerlessIsMore}.
Il modello di fatturazione \emph{pay-as-you-go} e la scalabilità
automatica hanno contribuito alla sua ampia adozione nell'industria.

All'interno dell'ecosistema serverless, il modello
\emph{Function-as-a-Service} (FaaS) rappresenta la forma
più diffusa di esecuzione event-driven.
In una piattaforma FaaS, l'unità fondamentale di
esecuzione è una funzione stateless, invocata in risposta a
eventi o richieste e tipicamente eseguita all'interno di
container o ambienti isolati equivalenti.

L'estensione del paradigma FaaS al continuum Edge--Cloud
risulta particolarmente promettente: eseguire funzioni
direttamente su nodi edge consente di ridurre la latenza
di risposta e migliorare la qualità del servizio percepita.
Recenti lavori hanno mostrato come architetture FaaS
decentralizzate possano abilitare l'esecuzione distribuita
di funzioni tra edge e cloud, introducendo meccanismi di
offloading verticale e orizzontale \cite{Russo2023_Serverledge}.
